\documentclass{simpledoc}


\newcommand{\soft}{\textsc{d-Tracker}}
\newcommand{\var}[1]{{\color{myblue}$\langle$#1$\rangle$}}
\newcommand{\opt}[1]{\color{myblue}$[$#1$]$}
\newcommand{\file}[1]{\textit{\color{myblue}#1}}
\newcommand{\screen}[1]{\texttt{\color{myblue}#1}}
\newcommand{\scr}[1]{\texttt{\color{myblue}#1}}
\newcommand{\key}[1]{\fbox{\color{myblue}#1}}

\title{\textsc{Tracker} -- Documentation}
\author{
\textbf{Vincent Richefeu} $\langle$Vincent.Richefeu@3sr-grenoble.fr$\rangle$\\
\textbf{Ga\"{e}l Combe} $\langle$Gael.Combe@3sr-grenoble.fr$\rangle$\\
Laboratoire 3SR, Universit\'{e} Grenoble Alpes, France
}
\date{Version 2026}

\pagestyle{myfootings}
\markboth{Tracker, user guide}%
         {Tracker, user guide}

%%%%%%%%%%%%%%%%%%%%%%%%%%%%%%%%%%%
%%%%%%%%%%%%%%%%%%%%%%%%%%%%%%%%%%%
\begin{document}

\maketitle

\begin{abstract}
This document describes how to use \soft, a code dedicated to 2D digital image correlation (DIC).
It covers particle image tracking, lens distortion correction, subpixel center detection, speckle
pattern quality analysis, gray level analysis, and result visualization.
\end{abstract}

\tableofcontents

\newpage

% ============================================================
\section{Introduction}

\soft\ is dedicated to 2D digital image correlation (DIC). It is written in C++ (standard C++17)
and can be compiled on Linux, and Mac OS X (and probably on Windows without the graphic applications).
It is actually a re-writting of the application \textsc{Tracker} developped in Laboratoire 3SR since 2011.
the application relies on the library \screen{libraw} to read raw camera images (CR2, IIQ, NEF, etc.) and on
\screen{libtiff} to read and write TIFF images.
Parallel computation is supported via OpenMP.

To compile, type \screen{make} in the \file{src} folder.
The resulting binary is \screen{tracker}, that is the CLI application. Recently, another binary has been added: \screen{tracker-gui}. It helps to visualise and tune some parameters, and its main goal is to help for some corrections in already tracked motions.


\attention
{\color{red} This code is designed for academic research. To be used with your own responsability.}

\subsection{Usage}





\begin{verbatim}
$ tracker <command_file>
\end{verbatim}

where \scr{command\_file} is a plain-text file (conventionally with extension \file{.trk})
that sets all parameters for the computation.
Running \screen{tracker} without arguments starts the interactive mode (see Section~\ref{sec:interactive}).

\subsection{Command file format}


A command file is a sequence of \texttt{keyword [value(s)]} pairs, one per line.
Lines starting with \texttt{!} or \texttt{\#} are comments and are ignored.
Inline comments are also possible (but add at least one space before the comment symbol):

\begin{verbatim}
! This is a full-line comment
search_zone.left  6    ! This is an inline comment
\end{verbatim}

\noindent The order of keywords in the file does not matter, except for a few cases where
a keyword must appear after the data it depends on (e.g.\ \scr{imposedDisplApprox\_corrDisto}
must follow \scr{image\_numbers\_corrDisto}).

% ============================================================
\section{Available procedures}
\label{sec:procedures}

The keyword \scr{procedure} selects which computation to run.
If omitted, \scr{particle\_tracking} is used by default.

\begin{center}
\begin{tabular}{lp{5.5cm}}
\hline
\textbf{Keyword} & \textbf{Description} \\
\hline
\scr{particle\_tracking}        & Default. Track particles between images. \\
\scr{correction\_distortion}    & Identify lens distortion parameters (displacement method). \\
\scr{correction\_distortion\_grid} & Identify lens distortion parameters (grid method). \\
\scr{find\_subpixel\_centers}   & Refine particle center positions to sub-pixel accuracy. \\
\scr{pattern\_quality}          & Assess speckle pattern quality by auto-correlation. \\
\scr{gray\_level\_analysis}     & Monitor grey level statistics over an image sequence. \\
\scr{visu\_process}             & Generate annotated visualization images from DIC results. \\
\scr{post\_process}             & Post-process DIC files (undistort, rescale, rotate). \\
\hline
\end{tabular}
\end{center}

% ============================================================
\section{Common keywords}

These keywords are shared by most procedures.

\marginlabel{\scr{procedure} \var{name}}
Select the computation procedure (see Section~\ref{sec:procedures}).

\marginlabel{\scr{image\_name} \var{path-format}}
Path to the image files, in C \texttt{printf} style.\\
Example: \texttt{image\_name ./images/img\%04d.tif}

\marginlabel{\scr{dic\_name} \var{path-format}}
Path pattern for the DIC output files.\\
Example: \texttt{dic\_name ./dic\_out\_\%d.txt}

\marginlabel{\scr{iref} \var{n}}
Number of the reference image.

\marginlabel{\scr{ibeg} \var{n}}
Number of the first image to process.

\marginlabel{\scr{iend} \var{n}}
Number of the last image to process.

\marginlabel{\scr{iinc} \var{n}}
Step between image numbers. Default is 1.

\marginlabel{\scr{idelta} \var{n}}
Offset between the reference image number and the DIC file number (used by \scr{visu\_process}).

\marginlabel{\scr{wanted\_num\_threads} \var{n}}
Number of OpenMP threads to use. The maximum is set at compile time (default 48).

\marginlabel{\scr{RawImages} \opt{0$\mid$1}}
Set to 1 to read raw camera files with \textsc{libraw} instead of TIFF files.

\marginlabel{\scr{DemosaicModel} \var{n}}
Demosaic algorithm passed to \textsc{libraw} (0--11). Default is 0.
Common values: 0 = linear, 3 = VNG, 9 = LMMSE, 11 = DHT.

\marginlabel{\scr{rescaleGrayLevels} \opt{0$\mid$1}}
If set to 1, rescale gray levels to span the full 16-bit range after reading.

\marginlabel{\scr{verbose\_level} \var{n}}
Controls the amount of output printed to the console. Default is 0.

% ============================================================
\section{The Particle Image Tracking procedure}
\label{sec:pit}

The \scr{particle\_tracking} procedure tracks the displacement and rotation of a set of points
(``grains'') between a fixed reference image and a sequence of deformed images.
For each point, a small region of the reference image (the \emph{pattern}) is correlated against
the deformed image over a search zone.
The result is a displacement vector $(u, v)$ and an incremental rotation $\Delta\theta$,
refined to sub-pixel accuracy by Powell minimization of $(1-\mathrm{NCC})$.

\subsection{Points to follow}

Three methods are available to define the tracked points.

\marginlabel{\scr{grains\_to\_follow} \var{filename}}
(also accepted as \scr{start\_file})
Read an ASCII file containing the initial positions of the tracked points.
File format:
\begin{verbatim}
N
x1  y1  rot1  r1
x2  y2  rot2  r2
...
\end{verbatim}
where $N$ is the number of points, $(x_i, y_i)$ are pixel coordinates (integers) on the reference
image, $\mathrm{rot}_i$ is the initial orientation in radians, and $r_i$ is the radius in pixels.

\attention The $y$ axis points downward (top of the image is $y=0$).

By convention, corners are assigned $r=1$ and fixed reference points $r=2$;
these special radii are used by the post-processing step to distinguish
sample elements from boundary markers.

\marginlabel{\scr{make\_grid} \var{xmin} \var{xmax} \var{ymin} \var{ymax} \var{nx} \var{ny} \var{alea}}
Generate a regular rectangular grid of $nx \times ny$ tracking points spanning the pixel rectangle
$[xmin, xmax] \times [ymin, ymax]$.
The parameter \scr{\var{alea}} adds a random jitter of up to \scr{\var{alea}} pixels to each point position
(useful to avoid aliasing artefacts). Set to 0 for a perfectly regular grid.
Example: \texttt{make\_grid 100 1000 100 1000 10 10 5}

\marginlabel{\scr{make\_grid\_circular} \var{xc} \var{yc} \var{R} \var{da}}
Generate a circular ring of points centered at $(xc,\ yc)$, at radius $R$ pixels,
with angular spacing $d\alpha$ radians.

\marginlabel{\scr{restart\_file} \var{filename}}
Read a previously saved DIC output file as the initial state, continuing from where a
previous run ended.
Set \scr{ibeg} to the image number just after the restart file.

\marginlabel{\scr{reordering} \var{order}}
Reorder the loaded grain list and save it to a new file, then exit.
Accepted values:\\
\scr{increasingRadius}, \scr{decreasingRadius}, \scr{increasingHeight}, \scr{mixed}.

\subsection{Image flow}

\marginlabel{\scr{iref} \var{n}}
Number of the reference image. The program always correlates against this fixed image,
avoiding cumulative errors.

\marginlabel{\scr{ibeg} \var{n}}
Number of the first deformed image to process.

\marginlabel{\scr{iend} \var{n}}
Number of the last deformed image to process.

\marginlabel{\scr{iinc} \var{n}}
Step between successive image numbers. For example, \texttt{iinc 5} processes
images \scr{ibeg}, \scr{ibeg}$+5$, \scr{ibeg}$+10$, etc.

\marginlabel{\scr{iraz} \var{n}}
Periodicity for resetting the cumulated displacement to zero (not yet fully implemented;
set to a value larger than \scr{iend} to disable).

\marginlabel{\scr{idelta} \var{n}}
Offset between the image index and the DIC file index (see \scr{visu\_process}).

\subsection{Shape of the pattern}

The pattern is the region of the reference image used for correlation.
All pattern commands apply to every grain unless \scr{targetRadiusPattern} is set.

\marginlabel{\scr{pattern circ} \var{R}}
Circular (disk) pattern of radius $R$ pixels (integer).
Example: \texttt{pattern circ 10}

\marginlabel{\scr{pattern ring} \var{Rin} \var{Rout}}
Annular pattern between inner radius $Rin$ and outer radius $Rout$ (both integers in pixels).
Example: \texttt{pattern ring 15 25}

\marginlabel{\scr{pattern rect} \var{dx} \var{dy}}
Rectangular pattern of half-width $dx$ and half-height $dy$ (integers in pixels),
giving a window of size $(2\,dx+1) \times (2\,dy+1)$.
Example: \texttt{pattern rect 10 10}

\marginlabel{\scr{pattern grain\_diameter} \var{factor} \var{min\_radius}}
Circular pattern whose radius is the grain radius multiplied by \scr{\var{factor}},
with a minimum of \scr{\var{min\_radius}} pixels.
Useful when grains have different sizes.

\marginlabel{\scr{pattern file} \var{filename}}
Load a custom pattern from an ASCII file (list of relative $(dx, dy)$ offsets).

\marginlabel{\scr{targetRadiusPattern} \var{r}}
Apply the next \scr{pattern} command only to grains whose radius equals exactly \scr{\var{r}}.
Set to a negative value (default) to apply to all grains.

\marginlabel{\scr{mask\_rect} \var{xmin} \var{xmax} \var{ymin} \var{ymax}}
Disable (mask) all grains whose reference position falls inside the rectangle
$[xmin, xmax] \times [ymin, ymax]$.
Masked grains are kept in the output file but not tracked.

\subsection{Search zone and rescue levels}

\marginlabel{\scr{rescue\_level} \opt{0$\mid$1$\mid$2}}
Controls how aggressively the program searches for poorly-matched grains:
\begin{itemize}
\item 0: no rescue (use the result of the first search regardless of quality)
\item 1: activate rescue with an extended search zone if $\mathrm{NCC} < \scr{NCC\_min}$
\item 2: additionally activate super-rescue if $\mathrm{NCC} < \scr{NCC\_min\_super}$ after rescue
\end{itemize}

\marginlabel{\scr{rotations} \opt{0$\mid$1}}
Whether to test rotations during the pixel-resolution search. Set to 0 for translating-only
patterns (faster). Default is 1.

The search zone parameters below control the pixel-resolution phase (before sub-pixel refinement).
Separate zones are defined for the first attempt, the rescue, and the super-rescue.

\marginlabel{\scr{search\_zone.left} \var{n}} Maximum displacement in pixels towards the left.
\marginlabel{\scr{search\_zone.right} \var{n}} Maximum displacement in pixels towards the right.
\marginlabel{\scr{search\_zone.up} \var{n}} Maximum displacement in pixels upwards.
\marginlabel{\scr{search\_zone.down} \var{n}} Maximum displacement in pixels downwards.
\marginlabel{\scr{search\_zone.inc\_rot} \var{$\Delta\theta$}}
Rotation step in radians. A good rule of thumb: $\Delta\theta \leq 1/R_{\max}$ rad,
where $R_{\max}$ is the distance (in pixels) from the center to the furthest pattern pixel.
\marginlabel{\scr{search\_zone.num\_rot} \var{n}}
Number of rotation increments tested in each direction;
$2n+1$ angles are tested in total (from $-n\,\Delta\theta$ to $+n\,\Delta\theta$).

\marginlabel{\scr{NCC\_min} \var{threshold}}
If the best NCC found in the first search is below this value, the rescue procedure is invoked.
Typical value: 0.8.

\marginlabel{\small \scr{search\_zone\_rescue.\opt{X} \var{?}}}
Same parameters as above but for the rescue search zone (should be larger).

\marginlabel{\scr{NCC\_min\_super} \var{threshold}}
If the best NCC after rescue is still below this value, super-rescue is invoked.
Typical value: 0.5--0.7.

\marginlabel{\small \scr{search\_zone\_super\_rescue.\opt{X} \var{?}}}
Same parameters but for the super-rescue search zone (should be even larger).

\marginlabel{\scr{use\_neighbour\_list} \opt{0$\mid$1}}
Use a neighbour list to guide the initial displacement guess for each grain
from the displacements of its neighbours. Improves robustness. Default is 1.

\marginlabel{\scr{num\_neighbour\_max} \var{n}}
Maximum number of neighbours stored per grain. Default is 40.

\marginlabel{\scr{neighbour\_dist\_pix} \var{d}}
Maximum distance in pixels to be considered a neighbour. Default is 250.

\marginlabel{\scr{period\_rebuild\_neighbour\_list} \var{n}}
Rebuild the neighbour list every \var{n} images. Default is 50.

\subsection{Sub-pixel refinement}

After the pixel-resolution search, a sub-pixel refinement is performed by minimizing
$(1 - \mathrm{NCC})$ using the Powell method with bicubic (or other) interpolation.

\marginlabel{\scr{subpixel} \opt{0$\mid$1}}
Activate (1) or deactivate (0) sub-pixel refinement. Default is 1.

\marginlabel{\scr{image\_interpolator} \opt{linear$\mid$cubic$\mid$quintic}}
Interpolation scheme used for sub-pixel image evaluation.
\scr{cubic} is the default and offers a good balance of accuracy and speed.
\scr{quintic} is more accurate but slower.

\marginlabel{\scr{subpix\_tol} \var{tol}}
Convergence tolerance for the Powell minimization. Default is $10^{-8}$.

\marginlabel{\scr{initial\_direction\_dx} \var{$\delta x$}}
\marginlabel{\scr{initial\_direction\_dy} \var{$\delta y$}}
\marginlabel{\scr{initial\_direction\_dr} \var{$\delta r$}}
Initial step sizes for the Powell minimizer in $x$, $y$ (pixels), and rotation (radians).
Should be small enough to converge to the global minimum. Default: $0.01$.

\marginlabel{\scr{subpix\_displacement\_max} \var{d}}
Maximum sub-pixel displacement allowed from the pixel-resolution result (in pixels).
Displacements exceeding this threshold are penalized by a stiffness term. Default is 3.0.

\marginlabel{\scr{overflow\_stiffness} \var{k}}
Stiffness of the penalty term for displacements exceeding \scr{subpix\_displacement\_max}.
Set to 0.0 to disable (default).

\subsection{Image generation for visual check}
\label{sec:imgcheck}

\soft\ can produce annotated images showing the tracked positions and correlation quality.

\marginlabel{\scr{make\_images} \opt{0$\mid$1}}
Activate (1) or disable (0) the generation of check images. Default is 0.

\marginlabel{\scr{image\_div} \var{n}}
Divide image dimensions by \var{n} when creating check images
(useful for large images). Default is 1 (full size).

\marginlabel{\scr{draw\_angle} \opt{0$\mid$1}}
Draw the rotation of each grain as a line. Default is 1.

\marginlabel{\scr{draw\_disp} \opt{0$\mid$1}}
Draw the displacement vector of each grain. Default is 1.

\marginlabel{\scr{draw\_rescued} \opt{0$\mid$1}}
Color code grains according to rescue status: green if rescue succeeded, red if it failed. Default is 1.

\subsection{Output format}
\label{sec:output}

For each processed image $i$, \soft\ writes a DIC file named\\
\file{dic\_out\_i.txt}.
The file begins with the total number of tracked points $N$, followed by one line per point:

\begin{center}
\begin{tabular}{rlp{7cm}}
\textbf{Col.} & \textbf{Variable} & \textbf{Description} \\
\hline
1 & \scr{refcoord\_xpix} & $x$ coordinate on the reference image (pixels) \\
2 & \scr{refcoord\_ypix} & $y$ coordinate on the reference image (pixels) \\
3 & \scr{refrot} & Reference orientation (radians) \\
4 & \scr{radius\_pix} & Grain radius (pixels) \\
5 & \scr{dx} & Cumulated $x$ displacement from reference (pixels) \\
6 & \scr{dy} & Cumulated $y$ displacement from reference (pixels) \\
7 & \scr{drot} & Cumulated rotation from reference (radians) \\
8 & \scr{upix} & $x$ displacement increment from last image (pixels) \\
9 & \scr{vpix} & $y$ displacement increment from last image (pixels) \\
10 & \scr{rot\_inc} & Rotation increment from last image (radians) \\
11 & \scr{NCC} & Best NCC from first search \\
12 & \scr{NCC\_rescue} & Best NCC after rescue (0 if no rescue was triggered) \\
13 & \scr{NCC\_subpix} & NCC after sub-pixel refinement (0 if no sub-pixel) \\
\hline
\end{tabular}
\end{center}

The file ends with two comment lines recording the image file names and timestamps.

If \scr{restart\_file} is used, the \emph{simple format} is also recognized:
$N$ lines containing only $(x, y, \mathrm{rot}, r)$ -- this is also the format produced by\\
\scr{grains\_to\_follow}.

% ============================================================
\section{Synthetic image generation}

Run \soft\ with the \texttt{-g} flag to enter a guided mode that generates a pair of synthetic
images using the Perlin noise algorithm.
You will be prompted for the following parameters:

\begin{itemize}
  \item \scr{width}, \scr{height}: image dimensions in pixels.
  \item \scr{zoom}: spatial scale of the noise pattern (larger = coarser speckle).
  \item \scr{octaves}: number of noise octaves (controls the richness of the texture).
  \item \scr{persistence}: weight of successive octaves (between 0 and 1).
  \item \scr{defx}, \scr{defy}: amplitude of the simulated deformation field.
  \item \scr{transx\_pix}, \scr{transy\_pix}: rigid-body translation (pixels).
  \item \scr{rot\_deg}: rigid-body rotation (degrees).
\end{itemize}

Two TIFF images are written: the reference image and the deformed image.
These can be used to validate the DIC algorithm and to study sensitivity to pattern quality.

% ============================================================
\section{Correction of distortion}
\label{sec:disto}

Lens distortion is modelled by Brown's formula.
Let $(x_u, y_u)$ be the undistorted coordinates and $(x_d, y_d)$ the distorted pixel coordinates.
With $\Delta x = x_u - X_c$, $\Delta y = y_u - Y_c$, $r^2 = \Delta x^2 + \Delta y^2$:

\begin{align}
\nonumber
x_d &= X_c + \Delta x\,(1 + K_1 r^2 + K_2 r^4 + K_3 r^6) \\ 
    &+
       (P_1(r^2 + 2\Delta x^2) + 2P_2\,\Delta x\,\Delta y)(1 + P_3 r^2) \\
\nonumber
y_d &= Y_c + \Delta y\,(1 + K_1 r^2 + K_2 r^4 + K_3 r^6) \\
    &+
       (P_1(r^2 + 2\Delta y^2) + 2P_2\,\Delta x\,\Delta y)(1 + P_3 r^2)
\end{align}

where $(X_c, Y_c)$ is the centre of distortion, $K_1, K_2, K_3$ are radial distortion
coefficients and $P_1, P_2, P_3$ are tangential distortion coefficients.

The optimization minimizes the error over the 8 parameters using Powell's method.
Initial guesses and perturbation amplitudes must be provided by the user.
The perturbation amplitude for a parameter can be set to zero to keep it fixed.

\subsection{Distortion correction parameters (both methods)}

\marginlabel{\scr{xc\_corrDistor} \var{$X_c$}}
$x$ coordinate of the distortion centre. A negative value is interpreted as half the image width.

\marginlabel{\scr{yc\_corrDistor} \var{$Y_c$}}
$y$ coordinate of the distortion centre. A negative value is interpreted as half the image height.

\marginlabel{\scr{\opt{K.}\_corrDistor} \var{$K$.}}
Radial distortion coefficients (initial guess). Typical starting value: 0.


\marginlabel{\scr{\opt{P.}\_corrDistor} \var{$P$.}}
Tangential distortion coefficients (initial guess). Typical starting value: 0.

\marginlabel{\scr{d\opt{.}c\_corrDistor} \var{.}}
Initial perturbation amplitudes for the Powell optimizer.
Set a perturbation to 0 to hold the corresponding parameter fixed.

\marginlabel{\scr{fake\_undistor} \opt{0$\mid$1}}
If set to 1, uses a simplified (linearized) undistortion function instead of the full
iterative Newton inversion. This can converge faster and may be sufficient in practice.

\marginlabel{\scr{undistor} \var{filename}}
Apply the current distortion parameters to undistort a single image file.
The result is written to \file{undist\_filename}.

\subsection{Displacement method (\scr{procedure correction\_distortion})}

This method estimates the distortion parameters from several images of the same scene
taken at different known (approximate) translations.
The DIC is first run on the image pairs to measure the actual displacements,
then the distortion parameters are optimized so that the measured displacements are
\emph{equi-projective} (i.e., consistent with a single homography).

\marginlabel{\scr{image\_numbers\_corrDisto} \var{n} \var{$i_1$} \ldots \var{$i_n$}}
List of $n$ image numbers to use (including the reference image first).

\marginlabel{\small\scr{imposedDisplApprox\_corrDisto} \var{$u_2$ $v_2$} \ldots \var{$u_n$ $v_n$}}
Approximate displacement (in pixels) of each image relative to the first one.
One pair $(u, v)$ per image, starting with image 2.
The first image is implicitly $(0, 0)$.

\marginlabel{\scr{equiProj\_dist\_range} \var{dmin} \var{dmax}}
Only include in the optimization tracking points whose inter-point distance (in pixels)
falls between \var{dmin} and \var{dmax}.
This filters out points too close together (affected by local noise) or too far apart.

\marginlabel{\scr{setSolidSolutionFromRef} \var{dx} \var{dy} \var{drot}}
Manually initialize the displacement solution of all grains to the given rigid-body motion
before running the DIC. Useful to set approximate initial guesses.

\marginlabel{\small\scr{set\_solidTranslation} \var{u} \var{v}}
Shift the reference positions of all grains by $(u, v)$ pixels.
This is applied before the DIC, to compensate for a known translation.

\marginlabel{\small\scr{ShowSolidMotionError} \opt{0$\mid$1}}
If set to 1, compute and print the rigid-body motion error between the first and current image.

\subsection{Grid method (\scr{procedure correction\_distortion\_grid})}

This method uses a photograph of a regular calibration grid.
The positions of grid intersection points are provided manually (e.g.\ from ImageJ),
listed from top-left to bottom-right, row by row.
The distortion parameters are optimized so that the measured point positions align with
a perfect regular grid.

\marginlabel{\scr{grid\_image\_name} \var{filename}}
Path to the calibration grid image.

\marginlabel{\scr{start\_file} \var{filename}}
File containing the manually picked pixel positions of the grid intersection points,
listed from top-left to bottom-right\\
(same format as \scr{grains\_to\_follow}).

\marginlabel{\scr{grid\_dim} \var{nx} \var{ny}}
Number of grid intersection points in the $x$ and $y$ directions.
\textbf{This keyword is mandatory for the grid method.}

% ============================================================
\section{Subpixel definition of particle centers}
\label{sec:subpixcenters}

The procedure \scr{find\_subpixel\_centers} refines the pixel-level center positions
of particles (e.g.\ circular rods) to sub-pixel accuracy.
It works by finding the position where the radial gradient of gray level is maximized,
using Powell minimization.

A prerequisite is to know an approximate threshold value separating the particles
from the background. The procedure first writes a file \file{profile.txt} containing
horizontal gray-level profiles at three heights ($y = $ dim$_y$/4, dim$_y$/2, 3\,dim$_y$/4)
to help the user choose the threshold.

\paragraph{Required keywords}
\scr{image\_name}, \scr{iref}, \scr{grains\_to\_follow} (approximate initial positions).

\marginlabel{\small\scr{subpix\_center\_threshold} \var{val}}
Gray level threshold separating the particle interior (below) from the background (above).
Choose this value from the \file{profile.txt} gray-level profiles.

\marginlabel{\scr{subpix\_center\_dr0} \var{r}}
Search radius around the approximate initial center (pixels). Default is 10.

\marginlabel{\scr{subpix\_center\_xstep} \var{dx}}
\marginlabel{\scr{subpix\_center\_ystep} \var{dy}}
\marginlabel{\scr{subpix\_center\_rstep} \var{dr}}
Powell step sizes in $x$, $y$ (pixels), and radius. Default is 0.1.

\paragraph{Output}
Two files are written:
\begin{itemize}
  \item \file{recentered.data}: refined positions in the same format as \scr{grains\_to\_follow}.
  \item \file{recentered.log.data}: same but with additional diagnostic information.
\end{itemize}

% ============================================================
\section{Characterization of speckle patterns}
\label{sec:patternquality}

The procedure \scr{pattern\_quality} computes the normalized cross-correlation (NCC) as a
function of sub-pixel displacement for each tracked point.
This is essentially the auto-correlation of the reference image at each pattern location.

For each point, the NCC is computed at a dense grid of offsets within a $\pm 10$ pixel window
and binned by distance from the origin.
The mean NCC as a function of distance is a proxy for pattern quality: a steep initial slope
(sharp auto-correlation peak) indicates a rich, well-textured pattern suitable for DIC.

The required keywords are:\\
\scr{image\_name}, \scr{iref}, \scr{grains\_to\_follow} or \scr{make\_grid}.

\marginlabel{\scr{halfPatternQual} \var{n}}
Half-size (in pixels) of the square region used for each auto-correlation evaluation.
Smaller values give faster results; larger values are more representative of the actual
correlation kernel. Default is 9.

\paragraph{Output}
\begin{itemize}
  \item \file{autocor.txt}: auto-correlation data for each point.
    Columns: distance, NCC, slope, population count.
    Points are separated by blank lines.
  \item \file{autocorMeanDev.txt}: mean and standard deviation of NCC over all points,
    as a function of distance. Columns: distance, mean NCC, standard deviation.
\end{itemize}

% ============================================================
\section{Gray level analysis}
\label{sec:grayanalysis}

The procedure \scr{gray\_level\_analysis} computes statistics of gray levels within each pattern
over an image sequence.
This is useful to detect illumination changes, saturation, or drift during a test.

The required keywords are:\\
\scr{image\_name}, \scr{iref}, \scr{ibeg}, \scr{iend}, \scr{iinc},
\scr{grains\_to\_follow} (or \scr{make\_grid}), and a \scr{pattern} shape.

\paragraph{Output}
One file per tracked point: \file{gray\_level0.txt}, \file{gray\_level1.txt}, etc.
Each line corresponds to one image and contains:
\begin{enumerate}
  \item image number
  \item mean gray level
  \item absolute deviation
  \item standard deviation
  \item variance
  \item skewness
  \item kurtosis
  \item minimum gray level
  \item maximum gray level
\end{enumerate}

% ============================================================
\section{Visualization}
\label{sec:visu}

The procedure \scr{visu\_process} generates annotated TIFF and PPM images from a
series of DIC output files.
Each output image shows the deformed image with the tracked quantities (displacement or NCC)
mapped as colors.

\paragraph{Required keywords}
\scr{image\_name}, \scr{dic\_name}, \scr{iref}, \scr{ibeg}, \scr{iend}, \scr{iinc}.

\marginlabel{\scr{visu\_mode} \opt{discrete$\mid$continuum}}
\scr{discrete}: draw a colored disk at each tracked point.
\scr{continuum}: build a Delaunay triangulation of the tracked points and draw a
colored continuum field (strain, displacement, NCC) over the triangles.

\marginlabel{\scr{visu\_output} \var{field}}
Quantity to visualize. Accepted values:
\begin{itemize}
  \item \scr{NCC}: normalized cross-correlation coefficient.
  \item \scr{disp}: displacement magnitude $\sqrt{dx^2 + dy^2}$.
  \item \scr{exx}, \scr{eyy}, \scr{exy}: strain components (continuum mode only).
\end{itemize}

\marginlabel{\scr{colorMin} \var{min}}
\marginlabel{\scr{colorMax} \var{max}}
Minimum and maximum values of the color scale.
Values below \scr{colorMin} are colored blue; values above \scr{colorMax} are colored red.

\marginlabel{\scr{visu\_autoscale} \opt{0$\mid$1}}
If set to 1, automatically adjust the color range to the data range. Default is 1.

\marginlabel{\scr{visu\_alpha} \var{$\alpha$}}
Transparency blending coefficient between the color overlay and the original image
($0$: overlay only, $1$: image only). Default is 0.3.

\marginlabel{\scr{nx\_grid\_visu} \var{n}}
In continuum mode, number of grid cells used for the interpolated color field.

\marginlabel{\scr{smoothDegreeF} \var{n}}
Smoothing degree applied to the field before visualization. Default is 1.

\marginlabel{\scr{visu\_draw\_in\_ref} \opt{0$\mid$1}}
If 1, draw the patterns at their reference positions rather than at the deformed positions.

\paragraph{Output}
For each image: \file{Visu\_\textit{field}\_\textit{num}.ppm} (color visualization)
and \file{Visu\_\textit{field}\_\textit{num}.tiff}.
A color scale is written to \file{ColorMap.ppm} on the first image.

% ============================================================
\section{Post-processing}
\label{sec:postpro}

The procedure \scr{post\_process} applies geometric corrections to the raw DIC output files
and produces \file{CONF\textit{xxxx}} configuration files used by downstream analysis tools.
The steps applied in sequence are:
\begin{enumerate}
  \item Undistortion of pixel coordinates using the previously identified lens distortion parameters.
  \item Rescaling from pixels to physical units using two known fixed points.
  \item Rigid-body rotation removal using the two fixed reference points.
\end{enumerate}

Each output file \file{CONF\textit{xxxx}} contains:
\begin{itemize}
  \item the number of sample points followed by their $(x, y, \mathrm{rot}, r)$ coordinates,
  \item the four corner points,
  \item the fixed reference points,
  \item the scale factor (length per pixel),
  \item the source image file names.
\end{itemize}

\paragraph{Required keywords}
\scr{ibeg}, \scr{iend}, \scr{iinc}.
The input is automatically read from files named \file{dic\_out\_\textit{n}.txt}.

\marginlabel{\scr{post\_undistor} \opt{0$\mid$1}}
Apply undistortion. Default is 1. Requires valid distortion parameters to be set.

\marginlabel{\scr{post\_rescale} \opt{0$\mid$1}}
Rescale coordinates to physical units. Default is 1.
Requires at least 2 fixed reference points (radius $= 2$) and \scr{scaling\_distance}.

\marginlabel{\scr{scaling\_distance} \var{d}}
Known physical distance between the two first fixed reference points (in the desired output unit,
e.g.\ mm). Used to convert from pixels.

\marginlabel{\scr{post\_rotation} \opt{0$\mid$1}}
Remove the solid-body rotation of the sample frame between images. Default is 1.

\marginlabel{\scr{post\_sync} \opt{0$\mid$1}}
Synchronize with external sensor data. Placeholder (not yet implemented).

\marginlabel{\scr{radius\_corners} \var{r}}
Radius value used to identify corner points in the DIC files. Default is 1.

\marginlabel{\scr{radius\_fixedPoints} \var{r}}
Radius value used to identify fixed reference points in the DIC files. Default is 2.

% ============================================================
\section{Interactive mode}
\label{sec:interactive}

Running \screen{tracker} without arguments starts an interactive shell.
Type \screen{help} to list available commands:

\begin{center}
\begin{tabular}{lp{7cm}}
\hline
\textbf{Command} & \textbf{Description} \\
\hline
\scr{read\_image}    & Load a TIFF image from a file path. \\
\scr{read\_raw\_image} & Load a raw camera image using \textsc{libraw}. \\
\scr{DemosaicModel}  & Set the demosaic algorithm index. \\
\scr{minmax}         & Print the minimum and maximum gray values of the loaded image. \\
\scr{histo}          & Compute the histogram and save it to \file{histo.txt}. \\
\scr{image\_check}   & Write the loaded image as seen by \soft\ to \file{image\_check.tiff}. \\
\scr{q} / \scr{quit} & Exit interactive mode. \\
\hline
\end{tabular}
\end{center}

\end{document}
